\section{ANALISIS KOMPETITOR}

\subsection{Analisis Kompetitor}

Berikut adalah perbandingan Savora dengan platform serupa yang sudah ada di pasar berdasarkan data aktual per 2025--2026 \cite{mobilefoodwaste2024review}\cite{kusumawati2025mobile}:

\begin{landscape}
\begin{table}[H]
\centering
\caption{Perbandingan Fitur dengan Kompetitor}
\label{tab:kompetitor}
\small
\begin{tabularx}{\linewidth}{lXXXXX}
\toprule
\textbf{Aspek / Fitur}
  & \textbf{Savora}
  & \textbf{Too Good To Go}
  & \textbf{Surplus Indonesia}
  & \textbf{Garda Pangan}
  & \textbf{Olio} \\
\midrule
Jenis platform
  & Marketplace food rescue UMKM (MVP)
  & Global marketplace food rescue
  & Clearance sale app (recommerce)
  & Food bank sosial / nirlaba
  & P2P sharing + solusi bisnis \\
\addlinespace
Tersedia di Indonesia
  & \checkmark\ (MVP)
  & \ding{55}\ (baru Jepang, 2026)
  & \checkmark\ (aktif)
  & \checkmark\ (Surabaya \& sekitar)
  & Terbatas (komunitas minimal) \\
\addlinespace
Fokus UMKM lokal kuliner
  & \checkmark
  & Terbatas (resto besar)
  & \checkmark\ (UMKM, horeka)
  & \checkmark\ (distribusi gratis)
  & \ding{55} \\
\addlinespace
Transparansi produk
  & Tinggi (foto, deskripsi, kelayakan)
  & Rendah (\textit{Surprise Bag})
  & Sedang (nama \& foto)
  & Tinggi (SOP seleksi ketat)
  & Sedang \\
\addlinespace
Food Trust Index / info kelayakan
  & \checkmark\ (rule-based, tampil ke user)
  & \ding{55}
  & \ding{55}
  & Seleksi internal (tidak ke user)
  & \ding{55} \\
\addlinespace
Metode pembayaran
  & COD + payment gateway (Midtrans)
  & Digital in-app (tanpa COD)
  & Digital saja (OVO, GoPay, DANA; tanpa COD)
  & Tidak ada (donasi)
  & Digital in-app \\
\addlinespace
Dynamic discount berbasis kelayakan
  & \checkmark\ (rule-based, 10--50\%)
  & \checkmark\ (AI, berbasis historis, sejak 2025)
  & \checkmark\ (hingga 80\%, tidak berbasis kelayakan)
  & \ding{55}
  & \ding{55} \\
\addlinespace
Pickup code \& batas waktu
  & \checkmark\ (pickup code + deadline)
  & Window waktu (tanpa pickup code)
  & Window waktu
  & \ding{55}
  & \ding{55} \\
\addlinespace
Waste Log untuk mitra UMKM
  & \checkmark
  & \ding{55}
  & \ding{55}
  & \ding{55}
  & \ding{55} \\
\addlinespace
Help Center food rescue
  & \checkmark\ (laporan terstruktur)
  & Ada (umum)
  & Tidak dipublikasikan khusus
  & \ding{55}
  & \ding{55} \\
\addlinespace
Impact tracking ke pengguna
  & \checkmark\ (estimatif, bukan klaim resmi)
  & \checkmark\ (CO\textsubscript{2}e per transaksi)
  & \checkmark\ (estimasi CO\textsubscript{2}/CH\textsubscript{4})
  & \ding{55}
  & \checkmark\ (3,67 kg CO\textsubscript{2}e per kg) \\
\addlinespace
Rating \& ulasan
  & \checkmark
  & \checkmark
  & \checkmark
  & \ding{55}
  & \checkmark \\
\addlinespace
Sertifikasi
  & ---
  & B-Corp
  & B-Corp (pertama di Indonesia)
  & --- (nirlaba)
  & B-Corp \\
\addlinespace
Skala (2025--2026)
  & MVP lomba
  & 120 juta user, 20+ negara
  & Ratusan ribu user, Jawa/Bali/Sulsel
  & Lokal Surabaya
  & 9 juta user, 62+ negara \\
\bottomrule
\end{tabularx}
\end{table}
\end{landscape}

\subsection{Analisis SWOT}

\begin{table}[H]
\centering
\caption{Analisis SWOT Savora}
\label{tab:swot}
\begin{tabularx}{\textwidth}{XX}
\toprule
\textbf{Strengths (Kekuatan)} & \textbf{Weaknesses (Kelemahan)} \\
\midrule
\begin{itemize}[leftmargin=*,nosep]
    \item Fitur lengkap terintegrasi (FES, Dynamic Pricing, Trust Index)
    \item Fokus pada UMKM lokal Indonesia
    \item Transparansi kelayakan makanan
    \item Mekanisme donasi terintegrasi
\end{itemize} \\ 
& \begin{itemize}[leftmargin=*,nosep]
    \item Butuh waktu untuk adopsi awal
    \item Bergantung pada jumlah UMKM mitra
    \item Membutuhkan data historis yang cukup
\end{itemize} \\
\addlinespace
\textbf{Opportunities (Peluang)} & \textbf{Threats (Ancaman)} \\
\midrule
\begin{itemize}[leftmargin=*,nosep]
    \item Kebijakan pengurangan food waste nasional
    \item Peningkatan kesadaran lingkungan
    \item Potensi ekspansi ke kota lain
    \item Kolaborasi dengan pemerintah daerah
\end{itemize} \\ 
& \begin{itemize}[leftmargin=*,nosep]
    \item Masuknya pemain besar dari luar negeri
    \item Perubahan regulasi
    \item Resistensi UMKM terhadap teknologi baru
\end{itemize} \\
\bottomrule
\end{tabularx}
\end{table}

\subsection{Model Bisnis}

Savora mengadopsi model bisnis \textit{freemium} dengan beberapa sumber pendapatan:

\subsubsection{Revenue Streams}
\begin{enumerate}[leftmargin=*]
    \item \textbf{Komisi Transaksi} — Sebesar 5-10\% dari setiap transaksi penjualan makanan surplus melalui platform.
    \item \textbf{Langganan Premium} — Paket premium untuk UMKM dengan fitur analitik lanjutan, prediksi produksi, dan laporan detail seharga Rp 99.000/bulan.
    \item \textbf{Iklan dan Promosi} — Layanan promosi produk untuk UMKM mitra yang ingin meningkatkan visibilitas.
    \item \textbf{Kerjasama Korporat} — Program CSR perusahaan yang ingin berkontribusi pada pengurangan \textit{food waste}.
\end{enumerate}

\subsubsection{Pricing Strategy}
\begin{table}[H]
\centering
\caption{Strategi Harga Savora}
\label{tab:pricing}
\begin{tabularx}{\textwidth}{lXX}
\toprule
\textbf{Paket} & \textbf{Harga} & \textbf{Fitur} \\
\midrule
Basic & Gratis & Fitur dasar, maks 50 produk/bulan \\
Pro & Rp 99.000/bulan & Analitik lanjutan, laporan detail, unlimited produk \\
Enterprise & Custom & Khusus jaringan UMKM besar, dedicated support \\
\bottomrule
\end{tabularx}
\end{table}

\subsection{Rencana Sustainability}

\begin{enumerate}[leftmargin=*]
    \item \textbf{Tahun Pertama} — Fokus pada akuisisi UMKM dan pengguna di Kota Bandung. Target mencapai \textit{break even} dari komisi transaksi dalam 12 bulan.
    
    \item \textbf{Tahun Kedua} — Ekspansi ke 3-5 kota besar lainnya (Jakarta, Surabaya, Yogyakarta). Meluncurkan paket premium dan program CSR.
    
    \item \textbf{Tahun Ketiga} — Menjadi platform referensi manajemen \textit{food waste} di Indonesia. Potensi pendanaan dari investor untuk percepatan pertumbuhan.
    
    \item \textbf{Jangka Panjang} — Ekspansi ke Asia Tenggara dan pengembangan produk turunan seperti pelatihan UMKM dan konsultasi \textit{sustainability}.
\end{enumerate}
